%%% FIELD proof-response-cost-geometry
Write \(d=|S_c|\), index the columns by \(t\in\mathcal T_c\), and let \(B\) be the \(d\times2^d\) matrix with columns \(b_t\).  Put \(a_t=\theta_t^2\).  The assumptions give
\[
c^\top v=\frac12\sum_{a\in S_c}|c_a|=1.
\]
The columns span \(\mathbb R^{S_c}\): if \(e_a\) is a coordinate vector, then
\[
e_a=b_{e_a}+b_{\mathbf1},\qquad -e_a=b_{\mathbf1-e_a}+b_{\mathbf0}.
\]
Thus every signed coordinate expansion can be rewritten with nonnegative coefficients, so \(B\lambda=\xi\), \(\lambda\ge0\), is feasible for every \(\xi\).  Its objective is nonnegative and finite on a feasible point.

For completeness, we derive attainment and duality rather than assuming them.  Consider the finitely generated cone
\[
\mathcal C=\operatorname{cone}\bigl(\{(b_t,a_t):t\in\mathcal T_c\}\cup\{(0,1)\}\bigr)
 \subset\mathbb R^{d+1}.
\]
This cone is closed.  Indeed, by repeatedly deleting a linear dependence from a conic representation while preserving nonnegativity, every element has a representation using at most \(d+1\) linearly independent generators.  Along a convergent sequence one may pass to one of the finitely many fixed generator sets; the coefficients then converge because the inverse on the span of that linearly independent set is continuous.  If \(p=V_c(\xi)\), feasible sequences show that \((\xi,p)\) is in the closure of \(\mathcal C\), hence in \(\mathcal C\).  In a representation
\[
(\xi,p)=\sum_t\lambda_t(b_t,a_t)+r(0,1),\qquad \lambda_t,r\ge0,
\]
one has \(\sum_ta_t\lambda_t\le p\); the definition of \(p\) forces equality and \(r=0\).  This proves primal attainment.

Weak duality follows immediately: if \(|\ell^\top b_t|\le a_t\), then for every primal-feasible \(\lambda\),
\[
\ell^\top\xi=\sum_t\lambda_t\ell^\top b_t
 \le \sum_t\lambda_ta_t.
\]
For the reverse inequality, fix \(q<p\).  The point \((\xi,q)\) is outside the closed convex cone \(\mathcal C\).  The elementary closest-point separation argument in Euclidean space supplies \((z,\beta)\ne0\) such that
\[
z^\top b_t+\beta a_t\ge0,\qquad \beta\ge0,\qquad
z^\top\xi+\beta q<0.
\]
Here \(\beta>0\): if \(\beta=0\), the inequalities for a type and its complement give \(z^\top b_t=0\) for every \(t\), and the spanning calculation above gives \(z=0\), a contradiction.  Set \(\ell=-z/\beta\).  Complementary types have \(b_{\mathbf1-t}=-b_t\) and \(a_{\mathbf1-t}=a_t\), whence
\[
|\ell^\top b_t|\le a_t,\qquad \ell^\top\xi>q.
\]
Letting \(q\uparrow p\), compactness of the dual feasible set, proved below, gives a limit attaining \(p\).  This establishes the displayed dual representation.

The type \(\mathbf0\) has \(a_{\mathbf0}=0\) and \(b_{\mathbf0}=-\mathbf1/2\), so every dual-feasible \(\ell\) satisfies \(\ell^\top\mathbf1=0\).  For the singleton type \(e_a\), this identity gives \(\ell^\top b_{e_a}=\ell_a\), hence \(|\ell_a|\le c_a^2\).  The dual set is therefore closed and bounded, hence compact.  It follows in particular that the maximum is attained; successively minimizing coordinates over compact optimizer faces also makes the stipulated lexicographically least optimizer well defined.

Positive homogeneity for nonnegative scalars and subadditivity follow by scaling and adding primal feasible vectors.  Evenness follows by replacing every coefficient at \(t\) by the same coefficient at \(\mathbf1-t\).  Consequently \(V_c\) is convex.  The displayed coordinate representations give a finite constant \(C_c\) such that \(V_c(z)\le C_c\|z\|_1\).  Subadditivity and evenness therefore yield
\[
|V_c(z)-V_c(z')|\le V_c(z-z')\le C_c\|z-z'\|_1,
\]
which proves finiteness and continuity.

The finite set \(\{t:\theta_t\ne0\}\) is nonempty because the type equal to one on the positive coordinates of \(c\) has score one.  Let
\(m_c=\min_{t:\theta_t\ne0}|\theta_t|>0\).  For every representation of \(v\),
\[
1=c^\top v=\sum_t\lambda_t\theta_t,
\qquad
\sum_t\lambda_t\theta_t^2
 \ge m_c\sum_t\lambda_t|\theta_t|
 \ge m_c.
\]
Thus \(\kappa(c)\ge m_c>0\).  If \(t^+_a=\mathbf1\{c_a>0\}\), then \(b_{t^+}=v\) and \(\theta_{t^+}=1\), so the single coefficient \(\lambda_{t^+}=1\) proves \(\kappa(c)\le1\).  Strong duality at \(v\) gives \((\ell^\star)^\top v=\kappa(c)\).  Dividing by the now-proved positive number \(\kappa(c)\), and using the zero-type and dual constraints, gives
\[
(g^\star)^\top v=1,\qquad (g^\star)^\top\mathbf1=0,\qquad
|(g^\star)^\top b_t|\le\theta_t^2/\kappa(c),
\]
as required.

%%% FIELD proof-exact-primitives
All calculations use the lexicographic order
\(000,001,010,011,100,101,110,111\).  Since \(\mathbf1^\top c^\dagger=0\), one has \(\theta_t=(c^\dagger)^\top t\), and direct substitution gives
\[
(\theta_t)_t=(0,-3/4,-1/4,-1,1,1/4,3/4,0).
\]
Set \(\ell=(5/8,-1/16,-9/16)\).  Its coordinates sum to zero, so \(\ell^\top b_t=\ell^\top t\), and hence
\[
((\ell^\top b_t)_t)=(0,-9/16,-1/16,-5/8,5/8,1/16,9/16,0).
\]
This vector is dual feasible because the absolute values are bounded respectively by
\((\theta_t^2)_t=(0,9/16,1/16,1,1,1/16,9/16,0)\).  Moreover,
\[
b_{000}+b_{101}+b_{110}=(1/2,-1/2,-1/2)=v
\]
and the corresponding primal cost is \(0+(1/4)^2+(3/4)^2=5/8\).  Weak duality and
\(\ell^\top v=5/8\) therefore prove \(\kappa(c^\dagger)=5/8\).  The optimizer is unique: the constraints at \(101\) and \(110\), together with \(\ell_1+\ell_2+\ell_3=0\), imply
\[
\ell_1=-\ell_2-\ell_3\le1/16+9/16=5/8;
\]
equality forces \(\ell_2=-1/16\) and \(\ell_3=-9/16\).  Thus it is in particular the selected lexicographically least optimizer, and division by \(5/8\) gives
\(g^\star=(1,-1/10,-9/10)\).  This proves the asserted eight pairs and all longitudinal identities.

For the transverse measures, their displayed coefficients sum to two.  Since a measure of mass two satisfies \(\sum_t\lambda_tb_t=\sum_t\lambda_tt-\mathbf1\), direct coordinate summation gives
\[
\sum_t\lambda_d^+(t)b_t=(0,3/4,-1/4)=d^\dagger,\qquad
\sum_t\lambda_d^-(t)b_t=-d^\dagger.
\]
Their costs are, in both cases,
\[
\frac34\left(\frac14\right)^2+\frac14\left(\frac34\right)^2
=\frac3{16},
\]
and \((c^\dagger)^\top d^\dagger=-3/16+3/16=0\).

The polynomial \(nX^3-1\) has derivative \(3nX^2\) and no common root with it, hence is square-free.  It is negative at zero, positive at two, and strictly increasing on the positive half-line, so it has exactly one positive root, \(\alpha_n=n^{-1/3}\); Sturm bisection provides a rational isolating interval.

On a sign cell, the masses of \(\lambda_v^{\sigma_x}\) and \(\lambda_d^{\sigma_h}\) are three and two.  Therefore \(\sum_ta_{\sigma,t}=0\), so \(\sum_tp_{n,x,h,t}^\dagger=1\).  All noncentral coordinates are nonnegative.  Each central coordinate is at least
\[
p_{0,t}^\dagger\{1-\alpha_n(3|x|+2|h|)\}\ge0
\]
under the stated condition; hence \(p_{n,x,h}^\dagger\) is a probability law.  The longitudinal measures have centered moments \(v\) and \(-v\), while the transverse measures have centered moments \(d^\dagger\) and \(-d^\dagger\).  Since the central law has centered mean zero, the sign-cell definition consequently gives
\[
\mathbb E[b_t\mid x,h]=\alpha_n(xv+hd^\dagger).
\]
Taking the inner product with \(c^\dagger\), using \((c^\dagger)^\top v=1\) and \((c^\dagger)^\top d^\dagger=0\), yields
\(\mathbb E[\theta_t\mid x,h]=\alpha_nx\).  Finally, a schedule with type counts \(m\) has
\[
\tau_{c^\dagger}(m)=\frac1n\sum_tm_t(c^\dagger)^\top t
=\frac1n\sum_tm_t\theta_t,
\]
which is the required finite-population target identity, not a prior-mean substitution.

For rational \(T\) and \(\varepsilon\), the factorial formula for \(P_m(x^{\mathrm{obs}}\mid r)\) and the last target formula are rational.  On each of the four sign cells, the multinomial mass is a polynomial in \(x,h,\alpha_n\) with rational coefficients, and \(w_T\rho_\varepsilon\) is also such a polynomial.  Exact integration over rational rectangles therefore places every prior weight in \(\mathbb Q(\alpha_n)\).  This field is closed under the finite sums, products, and the divisions by positive \(D_{r,x}\) used to form \(D_{r,x}\), \(N_{r,x}\), \(S_{r,x}^\dagger\), and \(\mathcal B_{n,\varepsilon,T}^\dagger(r)\).  A minimum over the finite set \(\mathcal R_{n,3}\) is one of those field elements.  Signs of elements of the fixed algebraic extension are decided by refining the isolating interval for \(\alpha_n\) (with the rational perfect-cube case treated in \(\mathbb Q\)); hence every branch and the finite minimum are exact.

%%% FIELD stmt-unit-airy-feedback
Let \(\phi\) be the positive even \(L^2(\mathbb R)\)-normalized ground state of \(-4d^2/dx^2+|x|\), let \(C_A\) be its eigenvalue, and put \(h=-2\phi'/\phi\).  Then \(h^2-2h'=|x|-C_A\).  The functions \(h'\) and \(h^2\) are globally Lipschitz, \(h'\) is bounded, and there is a finite constant \(D\) such that \(|x h'(x)|\le D\sqrt{|x|}\) for all \(x\in\mathbb R\).

%%% FIELD stmt-qstar-upper
Suppose \(K\ge3\) is fixed, \(\sum_ac_a=0\), \(\sum_a|c_a|=2\), and \(|S_c|\ge3\).  The estimator \(\delta_n^\star\) is total.  There is a finite constant \(B_c\), depending only on the fixed contrast, such that for every \(n\ge1\),
\[
\max_{y\in\{0,1\}^{n\times K}}\mathbb E\!\left[(\delta_n^\star-\tau_c(y))^2\right]
\le \frac1n-\Lambda_cn^{-4/3}+B_cn^{-3/2}.
\]
Consequently,
\[
\liminf_{n\to\infty}n^{4/3}\{n^{-1}-R_n(c)\}\ge\Lambda_c,
\qquad
\liminf_{n\to\infty}n^{4/3}\{n^{-1}-R_n^{\mathrm{ind}}(c)\}\ge\Lambda_c.
\]

%%% FIELD proof-qstar-upper
Only active arms have positive design probability, and on them \(q_a^\star=|c_a|/2>0\).  Thus the displayed ratios defining \(T_i\) and \(U_i\) are well defined on every positive-probability observation; returning zero on the complementary null set makes the rule total.  The normalization assumptions imply that every \(\tau_c(y)\) lies in \([-1,1]\), so metric projection onto \([-1,1]\) cannot increase squared loss.  It suffices to analyze the unprojected rule.

For a fixed schedule put
\[
\theta_i=c^\top(Y_i-\mathbf1/2)=c^\top Y_i,\qquad
\gamma_i=(g^\star)^\top(Y_i-\mathbf1/2).
\]
The second equality for \(\theta_i\) uses \(\mathbf1^\top c=0\).  Direct expectation over the independent arm label gives
\[
\mathbb ET_i=\theta_i,\qquad \mathbb EU_i=\gamma_i.
\]
Because \((Y_i(a)-1/2)^2=1/4\), \(q_a^\star=|c_a|/2\), and \(\sum_a|c_a|=2\),
\[
\mathbb ET_i^2=\sum_{a\in S_c}\frac{c_a^2}{4q_a^\star}=1.
\]
Likewise, using \((g^\star)^\top v=1\),
\[
\mathbb E(T_iU_i)=\sum_{a\in S_c}\frac{c_ag_a^\star}{4q_a^\star}
=\frac12\sum_{a\in S_c}\operatorname{sign}(c_a)g_a^\star
=(g^\star)^\top v=1.
\]
Hence
\[
\operatorname{Var}(T_i)=1-\theta_i^2,\qquad
\operatorname{Cov}(T_i,U_i)=1-\theta_i\gamma_i.
\]
The dual inequalities give \(|\gamma_i|\le\theta_i^2/\kappa(c)\).  In particular, since the contrast is fixed and the alphabet finite, \(T_i\) and \(U_i\) are bounded uniformly in \(i,n,y\).

Let \(\kappa=\kappa(c)>0\).  Scaling the unit-slope Airy ground state from the Airy feedback lemma by
\(\phi_\kappa(x)=\kappa^{1/6}\phi(\kappa^{1/3}x)\) shows that its eigenvalue is
\(C_A\kappa^{2/3}=\Lambda_c\).  Its feedback \(h_c(x)=\kappa^{1/3}h(\kappa^{1/3}x)\) obeys
\[
h_c(x)^2-2h_c'(x)=\kappa|x|-\Lambda_c.
\]
The cited bounds and this fixed rescaling supply finite constants, depending only on \(c\), such that \(h_c'\) and \(h_c^2\) are globally Lipschitz, \(h_c'\) is bounded, and
\[
|x h_c'(x)|\le C_c\sqrt{|x|}.
\]

Put
\[
X=\sum_i(T_i-\theta_i),\quad W=\sum_i(U_i-\gamma_i),\quad
Q=\sum_i\theta_i^2,\quad P=\sum_i\theta_i\gamma_i,
\quad r=Q/n,\quad u=n^{-2/3}\sum_i\gamma_i.
\]
Independence across units and boundedness give, uniformly over schedules,
\[
\mathbb EX^2\le n,\qquad \mathbb EW^2\le C_cn,\qquad
\mathbb EW^4\le C_cn^2.
\]
The fourth-moment bound follows by expanding the fourth power: terms containing a singleton centered factor vanish, leaving \(O(n)\) fourth moments and \(O(n^2)\) products of second moments.  Cauchy--Schwarz then gives
\[
\mathbb E(|X|W^2)\le(\mathbb EX^2)^{1/2}(\mathbb EW^4)^{1/2}le C_cn^{3/2}.
\]

Set \(a_n=n^{-2/3}\).  Taylor's formula with the Lipschitz derivative gives
\[
h_c(u+a_nW)=h_c(u)+a_nh_c'(u)W+R,\qquad
|R|\le C_ca_n^2W^2.
\]
Because \(\mathbb EX=0\) and
\(\mathbb E(XW)=\sum_i\operatorname{Cov}(T_i,U_i)=n-P\),
\[
\mathbb E\{Xh_c(u+a_nW)\}=a_nh_c'(u)(n-P)+O_c(a_n^2n^{3/2}).
\]
The global Lipschitz bound for \(h_c^2\), together with
\(\mathbb E|W|\le(\mathbb EW^2)^{1/2}\), also gives
\[
\mathbb Eh_c(u+a_nW)^2=h_c(u)^2+O_c(a_n n^{1/2}).
\]
Substitution into the exact squared-error expansion therefore yields
\[
\mathbb E(\delta_{n,\mathrm{unproj}}^\star-\tau_c(y))^2
=\frac1n-\frac rn+n^{-4/3}\left[h_c(u)^2-2h_c'(u)(1-P/n)\right]
+O_c(n^{-3/2}),
\]
uniformly in \(y\).  Here the first term uses
\(\mathbb EX^2=n-Q\), while the two displayed Taylor remainders are respectively
\(O_c((a_n/n)a_n^2n^{3/2})=O_c(n^{-3/2})\) and
\(O_c(a_n^3n^{1/2})=O_c(n^{-3/2})\).

The score inequality gives
\[
r\ge\kappa n^{-1/3}|u|,\qquad |P|/n\le r/\kappa.
\]
Indeed, both follow after summing \(|\gamma_i|\le\theta_i^2/\kappa\) and using \(|\theta_i|\le1\).  Write
\(d=r-\kappa n^{-1/3}|u|\ge0\).  The Riccati identity transforms the excess of the preceding risk over
\(n^{-1}-\Lambda_cn^{-4/3}\) into a quantity at most
\[
-\frac dn+\frac{2\|h_c'\|_\infty}{\kappa}d\,n^{-4/3}
+2|u h_c'(u)|n^{-5/3}+O_c(n^{-3/2}).
\]
For all sufficiently large \(n\), the first two terms have nonpositive sum.  Also \(|u|\le C_cn^{1/3}\), so the Airy bound gives
\(2|u h_c'(u)|n^{-5/3}=O_c(n^{-3/2})\).  This proves the asserted inequality for all sufficiently large \(n\), uniformly in the schedule.  Enlarging \(B_c\) over the finitely many remaining \(n\) proves it for every \(n\ge1\).

The independent procedure is feasible in the unrestricted game and its estimator is feasible under the fixed independent design.  Hence
\[
R_n(c)\le R_n^{\mathrm{ind}}(c)
\le\max_y\mathbb E[(\delta_n^\star-\tau_c(y))^2].
\]
Subtracting from \(n^{-1}\), multiplying by \(n^{4/3}\), and taking lower limits proves the two final claims.

%%% FIELD stmt-orbit-bayes
Assume \(c=c^\dagger\), rational \(T\ge1\), rational \(\varepsilon>0\), and \(\alpha_n(3T+2\varepsilon)\le1\).  Under the permutation-invariant schedule prior \(\nu_{n,\varepsilon,T}^\dagger\), every outcome-independent assignment law \(\mathcal D_n\) and every possibly randomized labeled-data estimator \(\widehat\tau\) have Bayes risk at least
\[
\sum_{r\in\mathcal R_{n,3}}\pi(r)\mathcal B_{n,\varepsilon,T}^\dagger(r)
\ge B_{n,\varepsilon,T}^\dagger,
\]
where \(\pi\) is the allocation-count law induced by \(\mathcal D_n\).  Conditional on \(r\), the likelihood of the arm-success counts is exactly \(P_m(X^{\mathrm{obs}}\mid r)\), and the posterior mean attains \(\mathcal B_{n,\varepsilon,T}^\dagger(r)\).

%%% FIELD proof-orbit-bayes
The positivity condition and the exact-primitives proposition make \(\nu_{n,\varepsilon,T}^\dagger\) a genuine distribution on type counts.  Conditional on \(m\), draw the labeled schedule uniformly from the \(n!/\prod_tm_t!\) schedules in its orbit.  The outcome-independent design is selected without observing this schedule, so its induced count \(r\) is independent of \(m\); denote its law by \(\pi\).

Fix an assignment vector having count \(r\).  For a table \(H\) of response types by assigned arms, the number of labeled schedules realizing \(H\) is \(\prod_ar_a!/\prod_{t,a}H_{t,a}!\).  Dividing by the orbit size and summing precisely over the tables with the prescribed arm-success counts gives
\[
\Pr_m(X^{\mathrm{obs}}=x^{\mathrm{obs}}\mid A)
=\frac{\prod_ar_a!}{n!}
\sum_{H\in\mathcal H(m,r,x^{\mathrm{obs}})}
\prod_t\frac{m_t!}{\prod_aH_{t,a}!}
=P_m(x^{\mathrm{obs}}\mid r).
\]
Thus, conditional on \(r\), the particular labeled assignment and the locations of successes within each arm are ancillary for \(m\); the posterior depends on the observed labeled data only through \(x^{\mathrm{obs}}\).

For any realized \((r,x^{\mathrm{obs}})\), square completion gives
\[
\mathbb E[(\widehat\tau-\tau_{c^\dagger}(m))^2\mid r,x^{\mathrm{obs}}]
=\operatorname{Var}(\tau_{c^\dagger}(m)\mid r,x^{\mathrm{obs}})
+\mathbb E[(\widehat\tau-\mathbb E[\tau_{c^\dagger}(m)\mid r,x^{\mathrm{obs}}])^2
 \mid r,x^{\mathrm{obs}}].
\]
This remains valid for estimator randomization after adjoining its independent random seed.  The second term is nonnegative and is zero for the posterior mean.  Multiplying the posterior variance by the predictive mass \(D_{r,x}\) and summing over \(x^{\mathrm{obs}}\) gives exactly
\[
\sum_m\nu_{n,\varepsilon,T}^\dagger(m)\tau_{c^\dagger}(m)^2
-\sum_{x^{\mathrm{obs}}}\frac{N_{r,x}^2}{D_{r,x}}
=\mathcal B_{n,\varepsilon,T}^\dagger(r),
\]
with zero-mass cells contributing zero by definition.  Averaging over \(r\sim\pi\) and then bounding an average below by its minimum proves the result.

%%% FIELD proof-directional-spectral
Put \(d=|S_c|\) and \(\kappa=\kappa(c)>0\), with positivity supplied by the response-cost geometry lemma.  The dual optimizer gives, for every \(z\in\mathbb R^{S_c}\),
\[
V_c(z)\ge (\ell^\star)^\top z=\kappa(g^\star)^\top z.
\]
Applying the same inequality to \(-z\) and using evenness yields
\[
V_c(z)\ge\kappa|(g^\star)^\top z|. \tag{1}
\]
Moreover \((g^\star)^\top v=1\).

First take \(\psi\in C_c^\infty(\mathbb R^{S_c})\).  Use the linear coordinates
\(z=xv+w\), where \(x=(g^\star)^\top z\) and \(w\in\ker((g^\star)^\top)\).  Their constant nonzero Jacobian is denoted by \(J_g\).  Along each fiber, \(\partial_v\psi(xv+w)=\partial_x\psi(xv+w)\).  By the definition of \(C_A\) and the change of variables \(s=\kappa^{1/3}x\), every one-dimensional \(f\) in the scalar form domain satisfies
\[
\int_{\mathbb R}\{4|f'(x)|^2+\kappa|x||f(x)|^2\}\,dx
\ge C_A\kappa^{2/3}\int_{\mathbb R}|f(x)|^2\,dx. \tag{2}
\]
Applying (2) on every fiber, using (1), and then using Fubini gives
\[
\mathcal Q_c(\psi)
\ge C_A\kappa^{2/3}\|\psi\|_2^2.
\]
The form in \(\mathrm{def:directional-form}\) is the closure of this core; lower semicontinuity therefore extends the inequality to its entire closed form domain.  This proves the lower bound.

For the reverse bound, use the different direct sum
\(z=xv+u\), \(u\in\ker(c^\top)\), which is valid because \(c^\top v=1\).  Let its constant Jacobian be \(J_c\).  Given \(\zeta>0\), the definition and scaling of \(C_A\) provide \(f\in C_c^\infty(\mathbb R)\), \(\|f\|_2=1\), such that
\[
\int\{4|f'|^2+\kappa|x|f^2\}\le C_A\kappa^{2/3}+\zeta.
\]
Choose \(\rho\in C_c^\infty(\ker(c^\top))\) with \(\|\rho\|_2=1\), and for \(a>0\) put
\(\rho_a(u)=a^{-(d-1)/2}\rho(u/a)\).  The normalized product
\[
\psi_a(xv+u)=J_c^{-1/2}f(x)\rho_a(u)
\]
is smooth and compactly supported.  Evenness, positive homogeneity, and subadditivity of \(V_c\) give
\[
V_c(xv+u)\le |x|V_c(v)+V_c(u)=\kappa|x|+V_c(u).
\]
There is no derivative penalty in the transverse \(u\) directions, and homogeneity gives
\[
\mathcal Q_c(\psi_a)
\le C_A\kappa^{2/3}+\zeta
+a\int_{\ker(c^\top)}V_c(w)|\rho(w)|^2\,dw.
\]
The last integral is finite because \(\rho\) has compact support and \(V_c\) is continuous.  Letting \(a\downarrow0\) and then \(\zeta\downarrow0\) proves the upper bound.  This concentrating sequence is valid whether or not the transverse potential has flat directions; it makes no assertion that a normalized full-dimensional minimizer exists.  Combining the bounds gives the claimed value \(\Lambda_c=C_A\kappa(c)^{2/3}\).

%%% FIELD stmt-algebraic-complexity
Let \(n\ge1\), rational \(T\ge1\), rational \(\varepsilon>0\), and rational \(\delta>0\) satisfy \(\alpha_n(3T+2\varepsilon)\le1\).  Let \(L\) be the total binary encoding length of \(n,T,\varepsilon,\delta\).  There are an absolute integer \(C\) and a deterministic Turing algorithm which returns rationals \(B^-\) and \(B^+\) satisfying
\[
B^-\le B_{n,\varepsilon,T}^\dagger\le B^+,
\qquad B^+-B^-\le\delta,
\]
using at most \((n+L+\lceil\log_2(1+\delta^{-1})\rceil+2)^C\) bit operations.  Thus the exact eight-type orbit evaluator has polynomial worst-case bit complexity in \(n\), the rational input lengths, and \(\log(1/\delta)\); no computability assumption on arbitrary real contrasts is used.

%%% FIELD proof-algebraic-complexity
The dimensions eight response types and three arms are fixed.  We give an explicit exhaustive algorithm and bound every loop and coefficient height.

First detect whether \(n\) is a perfect cube.  In that case \(\alpha_n\) is rational.  Otherwise \(X^3-1/n\) has no rational root and is irreducible, so \(\mathbb Q(\alpha_n)\) has degree three with basis \(1,\alpha_n,\alpha_n^2\) and reduction rule \(\alpha_n^3=1/n\).  Integer cube testing, rational-root testing, and Sturm isolation for this fixed-degree polynomial use a number of bit operations polynomial in \(L\).

The finite index sets have the explicit bounds
\[
|\mathcal M_{n,3}|={n+7\choose7}\le(n+1)^7,\qquad
|\mathcal R_{n,3}|={n+2\choose2}\le(n+1)^2,\qquad
\prod_{a=1}^3(r_a+1)\le(n+1)^3.
\]
For fixed \((m,r,x^{\mathrm{obs}})\), enumerate every \(8\times3\) table with entries in \(\{0,\ldots,n\}\) and retain exactly those satisfying the displayed linear constraints.  This uses at most \((n+1)^{24}\) candidates.  Consequently all likelihoods are computed by at most \((n+1)^{36}\) candidate-table visits over all \((m,r,x^{\mathrm{obs}})\).  Precomputed factorials have \(O(n\log(n+1))\) bits, so each factorial quotient is obtained by integer arithmetic of polynomial bit complexity.

For each \(m\), resolve \(|x|\) and \(|h|\) separately on the four sign cells.  The integrand defining its prior weight is then a bivariate polynomial of degree at most \(n+8\) with coefficients in \(\mathbb Q(\alpha_n)\).  Repeated convolution and collection by the reduction \(\alpha_n^3=1/n\) use polynomially many coefficient operations.  Integrating each monomial over rational endpoints computes \(\nu_{n,\varepsilon,T}^\dagger(m)\) exactly.  The numerator and denominator bit lengths produced by powers of rational \(T,\varepsilon\), binomial coefficients, factorials, integrations, and polynomially many additions are bounded by a polynomial in \(n+L\).  This conclusion remains true for the deliberately naive rule of multiplying denominators at every addition, because the number of additions is itself polynomial in \(n\).

Now enumerate \((r,x^{\mathrm{obs}},m)\), form \(D_{r,x}\) and \(N_{r,x}\), decide whether \(D_{r,x}=0\), and, on positive cells, form \(N_{r,x}^2/D_{r,x}\).  Arithmetic in degree at most three reduces to a fixed number of rational operations: inversion is obtained from the determinant and adjugate of the \(3\times3\) multiplication matrix.  Thus all \(\mathcal B_{n,\varepsilon,T}^\dagger(r)\) are produced with polynomially bounded operation count and coefficient height.

For sign tests and the minimum over \(r\), use the signed-subresultant comparison scheme of the algebraic-complexity handle.  Zero testing is coefficientwise in the irreducible basis.  For a nonzero element \(z=p_0+p_1\alpha_n+p_2\alpha_n^2\), clear rational denominators and let \(H\) bound the resulting integer coefficient lengths.  Its field norm is a nonzero rational obtained as the determinant of a \(3\times3\) integer multiplication matrix.  The numerator of that norm has absolute value at least one, while its denominator and the two conjugate factors have at most \(O(H+\log(n+1))\) bits.  Therefore
\[
|z|\ge2^{-O(H+\log(n+1))}.
\]
Refining the isolating interval for \(\alpha_n\) by that many bisections fixes the sign of \(z\).  Equivalently, the signed subresultant sequence has fixed length at most three and coefficients of polynomial bit length.  Thus every zero test, positivity branch, and comparison is decided in polynomial bit complexity.  There are only \((n+1)^2\) candidates in the final minimum.

The selected value is a degree-at-most-three real algebraic number in \([0,1]\).  Bisect this rational interval and decide the sign at each midpoint by the same subresultant procedure until the interval width is at most \(\delta\).  This takes
\(O(1+\log_2(1/\delta))\) comparisons after the exact value has been selected.  Multiplying the polynomial bounds for the finitely many loops, fixed-degree field operations, coefficient heights, and bisection yields an absolute exponent \(C\) with the claimed bound.  The procedure terminates and its final interval has the asserted order and width.

%%% FIELD partial-qstar
The proved uniform-risk upper theorem establishes the total finite-data procedure, the uniform remainder, and both lower-limit inequalities.  What remains is the reverse inequality
\[
\limsup_{n\to\infty}n^{4/3}\{n^{-1}-R_n^{\mathrm{ind}}(c)\}\le\Lambda_c,
\]
uniformly against all labeled estimators under the fixed independent design.

%%% FIELD partial-allocation-tube
The exact eight-type identities validate the tube whenever \(\alpha_n(3T+2\varepsilon)\le1\).  The proved finite orbit lemma gives, without any independence-across-unit restriction,
\[
\operatorname{BayesRisk}(\mathcal D_n,\widehat\tau;\nu_{n,\varepsilon,T}^\dagger)
\ge\sum_r\pi(r)\mathcal B_{n,\varepsilon,T}^\dagger(r)
\ge B_{n,\varepsilon,T}^\dagger.
\]
No claimed coefficient expansion, shell limit, or matching primal construction follows from these exact finite identities.

%%% FIELD partial-sharp-limit
The attaining-side theorem and \(\kappa(c^\dagger)=5/8\) prove
\[
\liminf_{n\to\infty}n^{4/3}\{n^{-1}-R_n(c^\dagger)\}
\ge C_A(5/8)^{2/3}>0.
\]
The orbit lemma supplies a valid finite Bayes lower bound for every design, but no proved expansion of that bound identifies its limit or matches the attaining coefficient.

%%% FIELD partial-effective
For every fixed admissible rational input, the polynomial-complexity theorem gives a terminating exact evaluator for \(B_{n,\varepsilon,T}^\dagger\) and a requested-width rational enclosure.  This is a same-input computation result.  It does not prove that the scheduled shell, coefficient, or tube-removal errors converge uniformly for all later sample sizes.
